\section{Уравнения Эйлера-Лагранжа для системы.\\
Нуль-лагранжиан. Детерминант как нуль-лагранжиан}
\label{sec:Chapter3} \index{Chapter3}

\subsection{Уравнения Эйлера-Лагранжа для системы}
\noindent$M^{m \times n}$ --- Пространство вещественных $m \times n$ матриц.\\
$U \subset \R^n$ --- ограниченное открытое множество с гладкой границей
$\partial U$. \\
Лагранжиан имеет вид:
$$L \colon M^{m \times n} \times \R^m \times \overline{U} \to \R, \quad
L(P, z, x) = L(p_1^1, \dots, p_n^m, z^1, \dots, z^m, x_1, \dots, x_n),$$ 
$$P\in M^{m \times n},
\quad z\in\R^m, 
\quad x\in U$$ 
Рассмотрим функционал:
\begin{equation}
    I[w] = \int_U L(Dw(x), w(x), x) \, dx,
    \quad Dw(x) = \begin{pmatrix}
        w^1_{x_1} & w^1_{x_2} & \dots & w^1_{x_n} \\
        w^2_{x_1} & w^2_{x_2} & \dots & w^2_{x_n} \\
        \vdots & \vdots & \ddots & \vdots \\
        w^m_{x_1} & w^m_{x_2} & \dots & w^m_{x_n}
    \end{pmatrix} %% = \left( w_{x_i}^k \right)_{k=1,\dots,m; \, i=1,\dots,n}
\end{equation}
где $w \colon \overline{U} \to \R^m$, $w(x) = (w_1(x), \dots, w_m(x))$
--- гладкая вектор-функция ($w \in C^\infty(\overline{U}, \R^m)$), 
минимизирующая функционал $I[w]$ на множестве функций с заданными 
значениями на границе области:
\begin{equation}
    \mathcal{M} = \{ w \in C^\infty(\overline{U}, \R^m) \mid w|_{\partial U} = g \},
    \quad g \colon \partial U \to \R^m
\end{equation}

Варьируя функционал $i(\tau) = I[w + \tau v]$ по направлениям 
$v(x) \in C_0^\infty(U, \R^m)$ и приравнивая производную к нулю, 
получаем систему уравнений Эйлера-Лагранжа:
\begin{equation}
    - \sum_{i=1}^n \frac{\partial}{\partial x_i} L_{p_i^k} (Dw, w, x) + L_{z^k}(Dw, w, x) = 0, \quad \forall k = 1, \dots, m.
\end{equation}

\textbf{Подробный вывод:}\\
Фиксируем произвольную вектор-функцию 
$v(x) = (v^1(x), \dots, v^m(x)) \in C_0^\infty(U, \R^m)$. 
Рассмотрим однопараметрическое семейство функций $w(x) + \tau v(x)$, 
где $\tau \in \R$. 
Определим функцию одной переменной $i(\tau)$: 
\begin{equation}
i(\tau) = I[w + \tau v] = 
  \int_U L\big(Dw(x) + \tau Dv(x), w(x) + \tau v(x), x\big) , dx.
\end{equation}

Поскольку при $\tau = 0$ функция $w(x)$ доставляет минимум функционала, 
точка $\tau = 0$ является точкой локального экстремума для функции $i(\tau)$. 
Следовательно, по необходимому условию экстремума: 
\begin{equation}
\frac{di(\tau)}{d\tau} \bigg|_{\tau=0} = 0.
\end{equation}
Используя правило дифференцирования сложной функции, 
продифференцируем подинтегральное выражение по $\tau$. 
\begin{equation}
\frac{di(\tau)}{d\tau} = 
  \int_U \sum_{k=1}^m \left\{ 
    \sum_{i=1}^n L_{p_i^k}(Dw, w, x) v_{x_i}^k + L_{z^k}(Dw, w, x) v^k 
  \right\} dx = 0
\end{equation}
Применим операцию интегрирования по частям к первому слагаемому 
под интегралом для каждой компоненты $v_{x_i}^k$:
\begin{equation}
    \int_U L_{p_i^k}(Dw, w, x) v_{x_i}^k \, dx = 
     -\int_U (L_{p_i^k}(Dw, w, x))_{x_i} v^k \, dx
\end{equation}
где $\nu$ --- вектор внешней нормали к границе области. 
Поскольку направление варьирования на границе зануляется 
($\forall k \colon v^k|_{\partial U} = 0$), поверхностный интеграл полностью 
обращается в ноль.\\
Подставляя преобразованное выражение обратно в общее интегральное равенство 
вариации, меняем порядок суммирования по индексам $i$ и $k$:
\begin{equation}
    \int_U \sum_{k=1}^m \left\{ - \sum_{i=1}^n (L_{p_i^k}(Dw, w, x))_{x_i} + L_{z^k}(Dw, w, x) \right\} v^k(x) \, dx = 0
\end{equation}

Так как компоненты вектор-функции $v^k(x)$ выбираются независимо 
друг от друга для каждого $k = 1, \dots, m$, мы можем последовательно 
применять лемму Дюбуа-Раймона к каждой компоненте отдельно. 
В силу непрерывности лагранжиана и его производных, выражение в фигурных 
скобках обязано тождественно зануляться в любой точке области $U$.

Таким образом, получаем систему уравнений Эйлера-Лагранжа:
\begin{equation}
    - \sum_{i=1}^n (L_{p_i^k} (Dw, w, x))_{x_i} + L_{z^k}(Dw, w, x) = 0, \quad \forall k = 1, \dots, m.
\end{equation}


\subsection{Нуль-лагранжиан}
\begin{definition}
    Функция $L(P, z, x)$ называется \textit{нуль-лагранжианом}, 
    если система уравнений Эйлера-Лагранжа
    \begin{equation}
        - \sum_{i=1}^n \frac{\partial}{\partial x_i} L_{p_i^k} (Dw, w, x) + 
        L_{z^k}(Dw, w, x) = 0 \quad (k = 1, \dots, m)
    \end{equation}
    выполняется тождественно для \textbf{любой} гладкой функции $w(x)$.
\end{definition}

\begin{theorem}
    Пусть $L(P, z, x)$ --- нуль-лагранжиан. 
    Если $u(x)$ и $\tilde{u}(x)$ --- гладкие функции, совпадающие на границе 
    области (то есть $u(x) = \tilde{u}(x)$ при $x \in \partial U$), 
    то значения функционала на них равны:
    \begin{equation}
        I[u] = I[\tilde{u}].
    \end{equation}
\end{theorem}

\begin{proof}
    Рассмотрим вспомогательную функцию одной переменной:
    \begin{equation}
        i(\tau) = I[\tau u + (1-\tau)\tilde{u}] 
        = \int_U L(\tau Du + (1-\tau)D\tilde{u}, \tau u + (1-\tau)\tilde{u}, x) \, dx.
    \end{equation}
    Продифференцируем её по параметру $\tau$:
\begin{equation}
\begin{split}
    \frac{di(\tau)}{d\tau} = \int_U \bigg\{ 
    &\sum_{i=1}^n \sum_{k=1}^m L_{p_i^k}(
      \tau Du + (1-\tau)D\tilde{u}, \tau u + (1-\tau)\tilde{u}, x
      )(u_{x_i}^k - \tilde{u}_{x_i}^k) + \\
    &\sum_{k=1}^m L_{z^k}(
      \tau Du + (1-\tau)D\tilde{u}, \tau u + (1-\tau)\tilde{u}, x
      ) (u^k - \tilde{u}^k)
    \bigg\} dx.
\end{split}
\end{equation}

Проинтегрируем первое слагаемое по частям. 
Поскольку $u^k - \tilde{u}^k = 0$ на границе $\partial U$, 
внеинтегральные члены обратятся в ноль:

\begin{equation}
\begin{split}
    \frac{di(\tau)}{d\tau} = \int_U \sum_{k=1}^m \bigg( 
    &- \sum_{i=1}^n \frac{\partial}{\partial x_i} L_{p_i^k}(
      \tau Du + (1-\tau)D\tilde{u}, \tau u + (1-\tau)\tilde{u}, x
      ) + \\
    &L_{z^k}(
      \tau Du + (1-\tau)D\tilde{u}, \tau u + (1-\tau)\tilde{u}, x
      ) 
    \bigg) (u^k - \tilde{u}^k) \, dx.
\end{split}
\end{equation}

Так как $L$ --- нуль-лагранжиан, выражение в больших скобках равно нулю 
для любой гладкой функции, в том числе для выпуклой комбинации 
$\tau u + (1-\tau)\tilde{u}$. Значит, $\frac{di(\tau)}{d\tau} = 0$.
Отсюда $i(\tau) \equiv \mathrm{const}$, и следовательно, $i(1) = i(0)$, 
то есть $I[u] = I[\tilde{u}]$.
\end{proof}

\subsection{Детерминант как нуль-лагранжиан}

Введем обозначения. Пусть $A$ --- квадратная $n \times n$ матрица. Матрица алгебраических дополнений обозначается как $\operatorname{cof} A$, где ее элемент $(\operatorname{cof} A)_i^k = (-1)^{i+k} M_i^k$ ($M_i^k$ --- минор, полученный вычеркиванием $k$-й строки и $i$-го столбца). 
Известно, что для дифференциала определителя выполняется соотношение:
\begin{equation}
    \frac{\partial \det P}{\partial p_i^k} = (\operatorname{cof} P)_i^k.
\end{equation}

\begin{lemma}[о строках с нулевой дивергенцией]
    Пусть $u \colon \R^n \to \R^n$ --- гладкая вектор-функция. Тогда дивергенция любой строки матрицы алгебраических дополнений якобиана равна нулю:
    \begin{equation}
        \sum_{i=1}^n \frac{\partial}{\partial x_i} (\operatorname{cof} Du)_i^k = 0, \quad \forall k = 1, \dots, n.
    \end{equation}
\end{lemma}
\begin{proof}
Для любой матрицы $P \in \R^{n \times n}$ справедливо:
\begin{equation}
    \det P \cdot \delta_j^k = \sum_{i=1}^n p_i^k (\operatorname{cof} P)_i^j
\end{equation}
где $\delta_j^k$ --- символ Кронекера. Также напомним известное свойство дифференцирования определителя по элементам матрицы:
\begin{equation}
    (\det P)_{p_i^k} = (\operatorname{cof} P)_i^k
\end{equation}

Положим $P = Du(x)$. Тогда компоненты матрицы представляют собой частные производные по координатам: $p_i^k = u^k_{x_i}$. Подставив их в исходное алгебраическое тождество, получаем выражение:
\begin{equation}
    \det(Du) \cdot \delta_j^k = \sum_{i=1}^n u^k_{x_i} (\operatorname{cof} Du)_i^j
\end{equation}

Продифференцируем это тождество по переменной $x_s$:
\begin{equation}
    (\det(Du))_{x_s} \cdot \delta_j^k = \sum_{i=1}^n \left( u^k_{x_i x_s} (\operatorname{cof} Du)_i^j + u^k_{x_i} ((\operatorname{cof} Du)_i^j)_{x_s} \right)
\end{equation}

Распишем полную производную определителя $(\det(Du))_{x_s}$ по правилу дифференцирования сложной функции, используя свойство $(\det P)_{p_i^k} = (\operatorname{cof} P)_i^k$:
\begin{equation}
    (\det(Du))_{x_s} = \sum_{m=1}^n \sum_{l=1}^n (\operatorname{cof} Du)_l^m u^m_{x_l x_s}
\end{equation}

Свернем теперь индексы, положив $k = j$ и просуммировав по нему. В силу 
сокращений и тождественного совпадения сумм, содержащих вторые смешанные 
производные $u^k_{x_i x_s}$, члены с ними взаимно уничтожаются из левой и
правой частей уравнения. Это приводит к соотношению:
\begin{equation}
    \sum_{i=1}^n u^k_{x_i} \left( \sum_{s=1}^n ((\operatorname{cof} Du)_i^j)_{x_s} \right) = 0
\end{equation}

Если $\det(Du) \neq 0$, то матрица Якоби $Du$ (элементами которой являются $u^k_{x_i}$) является невырожденной. Данное соотношение можно рассматривать как матричное уравнение вида $Du \cdot \vec{v} = 0$, где вектор $\vec{v}$ состоит из сумм, стоящих в скобках. Так как матрица обратима, однородная система имеет только единственное (тривиальное) решение $\vec{v} = 0$. Умножая обе части соотношения на обратную матрицу $(Du)^{-1}$, мы доказываем, что выражение в скобках (дивергенция) тождественно равно нулю:
\begin{equation}
    \sum_{i=1}^n ((\operatorname{cof} Du)_i^k)_{x_i} = 0, \quad \forall k = 1, \dots, n.
\end{equation}

Если $\det(Du) = 0$ (в некоторых точках области), применим метод малых возмущений. Рассмотрим вспомогательную гладкую функцию $u_\epsilon(x) = u(x) + \epsilon x$. Матрица Якоби для нее примет вид $Du_\epsilon = Du + \epsilon I$, где $I$ --- единичная матрица.

Определитель $\det(Du + \epsilon I)$ представляет собой многочлен степени $n$ относительно параметра $\epsilon$ (характеристический многочлен матрицы $-Du$). Такой многочлен имеет не более $n$ корней, а значит, для почти всех сколь угодно малых $\epsilon > 0$ выполняется условие:
\begin{equation}
    \det(Du + \epsilon I) \neq 0.
\end{equation}

Согласно доказанному первому случаю, для возмущенной функции $u_\epsilon(x)$ тождество выполняется строго:
\begin{equation}
    \sum_{i=1}^n ((\operatorname{cof} Du_\epsilon)_i^k)_{x_i} = 0.
\end{equation}

Элементы матрицы алгебраических дополнений $(\operatorname{cof} P)_i^k$ и их частные производные представляют собой многочлены от элементов матрицы $P$, поэтому они зависят от $P$ непрерывно. Устремляя параметр возмущения к нулю ($\epsilon \to 0$), мы можем совершить предельный переход под знаком производных. В пределе получаем, что тождество справедливо и для исходной функции $u(x)$.
\end{proof}

\begin{theorem}
    Лагранжиан $L(P) = \det P$ является нуль-лагранжианом.
\end{theorem}

\begin{proof}
    Функция $L$ зависит только от матрицы производных $P = Du$. Проверим для нее выполнение системы уравнений Эйлера-Лагранжа:
    \begin{equation}
        - \sum_{i=1}^n \frac{\partial}{\partial x_i} L_{p_i^k}(Du) = 0.
    \end{equation}
    Подставляя частные производные лагранжиана, получаем:
    \begin{equation}
        - \sum_{i=1}^n \frac{\partial}{\partial x_i} \left( \frac{\partial \det P}{\partial p_i^k} \middle|_{P = Du} \right) = - \sum_{i=1}^n \frac{\partial}{\partial x_i} (\operatorname{cof} Du)_i^k = 0.
    \end{equation}
    Последнее равенство справедливо тождественно для любой гладкой функции $u(x)$ в силу леммы о строках с нулевой дивергенцией. Таким образом, $L(P) = \det P$ --- нуль-лагранжиан.
\end{proof}

\textbf{Пример Больца}\\
Рассмотрим вариационную задачу:
\begin{equation}
    I[x(t)] = \int_0^1 \left[ \left(1 - \left(\dot{x}\right)^2\right)^2 + x^2(t) \right] dt \rightarrow \min, \quad x(0) = x(1) = 0.
\end{equation}

Так как подынтегральная функция является суммой квадратов, $I[x] \ge 0$. 

Построим минимизирующую последовательность «пилообразных» функций 
$x_n(t) = \int_0^1 \operatorname{sign}(\sin(23n\pi \tau)) \, d\tau$. 
На линейных участках $\dot{x}_n(t) = \pm 1$, поэтому:
\begin{equation}
    I[x_n] = \int_0^1 x_n^2(t) \, dt \xrightarrow{n \to \infty} 0.
\end{equation}

Последовательность равномерно сходится к нулю ($x_n(t) \rightrightarrows 0$), 
и значение функционала на ней стремится к нулю. 
Однако для предельной функции $x(t) \equiv 0$ ($\dot{x}(t) \equiv 0$) 
значение функционала строго положительно:
\begin{equation}
    I[0] = \int_0^1 \left[ (1 - 0)^2 + 0 \right] dt = 1.
\end{equation}

